\lesson{4}{Oct 11 2021 Mon (11:28:43)}{The Anti-Federalists}{Unit 2}

\subsubsection*{The Constitution—To Ratify or Not to Ratify?}

When the \textbf{Constitution} was sent to the state legislatures for
ratification, not everyone was eager to approve it.  At least nine of the
thirteen states had to ratify for it to take effect. While the convention
drafting the document was still in session, delegates began to separate into two
groups. They were:

\begin{itemize}
    \item Federalists \\
        People who supported the Constitution as it was and argued for immediate
        ratification became known as "federal men" or Federalists. James Madison
        from Virginia was a Federalist who eventually served as president.
        Alexander Hamilton from New York also favored the Constitution and went
        on to be the first U.S. Secretary of the Treasury, writing economic
        policies for George Washington's administration.
    \item Anti-Federalists \\
        Those who did not wish to ratify were called \textbf{"anti-federals"}, though the
        name does not mean that they were against federalism. In fact, the
        Anti-Federalists' name is ironic because they were against a strong
        central government. They favored a federal system where the states were
        supreme, similar to the Articles of Confederation. Most of them agreed,
        however, that a somewhat stronger central government was necessary. Not
        all Anti-Federalists were the same. Some wanted very little deviation
        from the Articles of Confederation and were staunchly against the new
        Constitution. Others were more moderate and willing to compromise, while
        some focused only on adding the Bill of Rights. Most of those unwilling
        to compromise, like Richard Henry Lee and Patrick Henry, did not attend
        the Philadelphia Convention, but had a strong voice when the time came
        for states to vote on ratification. Lee, from Virginia, was "president
        of Congress" for a year under the Articles of Confederation. He went on
        to serve as a senator to Congress from his state. Henry twice served as
        governor of Virginia and refused an invitation to the convention, saying
        he "smelt a rat in Philadelphia, tending toward the monarchy." He served
        in the Virginia legislature during the ratification debates and for
        several years more.
\end{itemize}

Delegates met and debated in secret at the Philadelphia Convention. Pens were
busy once the document went to the states, with prominent leaders around the
country writing for or against its ratification. The lack of a unified
Anti-Federalist voice in Anti-Federalist articles is probably one reason why the
Federalists and the Federalist Papers they created are remembered more today.
However, as you will see, compromises between the Federalists and
Anti-Federalists had a significant impact on the Constitution then and now.

\subsubsection*{What Was the Great Compromise?}

Both Federalist and Anti-Federalist delegates at the Philadelphia Convention
agreed that a stronger central government was necessary than what the Articles
of Confederation provided. They also agreed that neither a pure democracy nor an
aristocracy were desirable or sustainable for the new nation. A system of
representation for each state was necessary. The greatest issue to divide
delegates at the constitutional convention was the strength of the central
government and the relative power between states with varying populations.

\begin{itemize}
    \item Conflicting Plans \\
        \paragraph{Virginia Plan}

        Under James Madison's "Virginia Plan," the nation would have a powerful central
        government and a bicameral legislature whose members would choose the president
        and federal judges. Federalists from states with large populations favored this
        plan. However, delegates from smaller states such as New Jersey, whether
        Federalist or Anti-federalist, feared that both houses of the legislature would
        depend on population in the new plan, giving greater power to larger states.

        \paragraph{New Jersey Plan}

        New Jersey delegates then proposed an alternate plan where the states retained
        the "one state, one vote" system. Anti-federalists were more likely to support
        this plan, since it would also have left the national government largely
        dependent on the states as it was before. Delegates from larger states did not
        think it fair for the states to have equal representation in Congress because
        some states had significantly greater numbers of people to serve. Delegates from
        smaller states preferred equal representation so their state would have as
        strong a voice in national matters as the larger states.

    \item Great Compromise \\
        In what we now call the Great Compromise, the delegates agreed that the
        Constitution should create a bicameral legislature. Congress would have
        two houses—the Senate and the House of Representatives. Each state would
        have two senators chosen by the state legislatures, as preferred by
        delegates from smaller states. The House of Representatives would be
        composed of a number of representatives based on population and elected
        directly by the people in each state, as preferred by delegates from
        larger states. The Senate and the House of Representatives, though they
        do have some distinctions in duties, must work together to achieve their
        main job of approving legislation before it becomes law.
\end{itemize}

\subsubsection*{How Else Does the Constitution Reflect Compromise?}

Small states and large states each achieved some of what they wanted in the
Great Compromise. It was only one compromise of many at the convention.

The delegates also debated how to address slavery in the Constitution, though it
never directly names it. Southern states, with large numbers of enslaved people,
wanted to count slaves into their population for determining representation to
Congress but not for taxation purposes. Northern states wanted the opposite.
Greater representation would give a larger state more power in passing federal
laws but also greater taxes. Delegates from northern states argued that counting
slaves would not be fair for representation because slaves were not citizens
with voting and other rights. Yet all states counted women and children, whose
rights were limited at best, into population. Like the Great Compromise, the
delegates to the convention were divided here more along state loyalty lines
rather than as Federalist against Anti-Federalist.

Some delegates hoped that slavery would end gradually and did not want to give
any protection or incentive for the practice to continue in the Constitution.
They wanted to end the slave trade so no new people from overseas would be
enslaved in the states. Yet all knew that without a compromise on this issue,
Southern States would not ratify the Constitution. They agreed to the
Three-Fifths Compromise. Each state would be able to count three-fifths of its
total enslaved people as part of the state population for representation
purposes. Delegates also wrote into the Constitution that Congress could not
pass any laws outlawing the slave trade until 1808. Of course, both of these
agreements became obsolete once slavery ended.

\subsubsection*{What Happened When the Constitution Went to the States?}

On September 17, 1787, the compromises were complete, and a majority of
delegates approved and signed the Constitution. Yet, nine state legislatures had
to ratify it before the new government could begin. Each state began its own
chain of discussion and debate over the document. People who had attended the
convention, as well as those who did not, wrote arguments for and against the
Constitution’s ratification that reached a public audience through newspapers.

\subsubsection*{How Did a Bill of Rights Help Ratification?}

Many Anti-Federalists would never have approved the Constitution without the
promise of a bill of rights. Federalists like James Madison argued that the
existing language in the Constitution was enough to protect individual rights,
especially considering the state constitutions with their own lists of protected
rights.

Anti-Federalists argued that if some rights were worthy of protection within the
document, then all rights should have the same level of written protection.

It was not that Federalists were against protecting individual rights. Rather,
some worried that listing them could backfire. They argued that the government
might use the document as justification to restrict a right because it does not
appear in print. Some Anti-Federalists saw it as a simple way to solidify
limitations on the power of the national government. Out of the over 200
suggestions for amendments, or changes, to the Constitution to address the
issue, James Madison wrote 12. The states ratified 10 of the amendments, which
were then named the Bill of Rights.

\subsubsection*{When Did the States Ratify?}

While some say the Anti-Federalists "lost" the debate, without them the
Constitution would not be the same. They pushed the delegates at the
constitutional convention to clarify the proper role of government at each
level. They secured the promise of a Bill of Rights to increase protection for
individual rights under the new government. In fact, some historians believe
that without a guarantee of amendments to come, the Constitution may never have
been ratified.

Achievements of the Anti-Federalists:

\begin{itemize}
    \item Amending the Constitution with a Bill of Rights
    \item Stronger protection for states' rights
\end{itemize}

Achievements of the Federalists:

\begin{itemize}
    \item Ratification of the Constitution
    \item Strong central government with separation of powers
\end{itemize}

State ratification was not immediate. Each state government had its own schedule
for debating and voting on the Constitution. The process took longer in some
states than others. Note, too, that just because a state ratified does not mean
everyone agreed with the Constitution. Examine this map to see the dates of
ratification for each state and the regional majorities of the Federalists and
Anti-Federalists.

\newpage
